\chapter{The Sound Key}

\section{Vowels}
Each English vowel is written as a short mark on the preceding consonant and, for long vowels and glides, a vowel letter. Table~\ref{tab:vowels} shows the default correspondences using words from the page. Phonemes are named in ARPAbet, the notation of the pronouncing dictionary.

\begin{table}[ht]\centering\small
\begin{tabular}{@{}lllll@{}}
\toprule
ARPAbet & Mark & Letters & Example & Rendered\\
\midrule
AA & fatha & \ar{ا} & part & \ar{بَارْتْ}\\
AE & fatha & -- & lack & \ar{لَاكْ}\\
AH & fatha & -- & cultural & \ar{كَلْتْشَرَلْ}\\
AO & damma & \ar{و} & on & \ar{أُونْ}\\
AW & fatha & \ar{او} & flower & \ar{فْلَاوَرْ}\\
AY & fatha & \ar{اي} & psychocinema & \ar{سَايْكُوسِينِمَا}\\
EH & kasra & -- & delves & \ar{دِلْفْزْ}\\
ER & fatha & \ar{ر} & nature & \ar{نَيْتْشَرْ}\\
EY & fatha & \ar{ي} & nature & \ar{نَيْتْشَرْ}\\
IH & kasra & -- & this & \ar{ذِسْ}\\
IY & kasra & \ar{ي} & she & \ar{شِي}\\
OW & damma & \ar{و} & Rollins & \ar{رُولِنْزْ}\\
OY & damma & \ar{وي} & void & \ar{فُويْدْ}\\
UH & damma & -- & could & \ar{كُودْ}\\
UW & damma & \ar{و} & into & \ar{إِنْتُو}\\
\bottomrule
\end{tabular}
\caption{Vowel key. The mark is the short vowel; the letters are the vowel letters added for long vowels and glides.}\label{tab:vowels}
\end{table}

Three kinds of vowel letter occur. A \emph{long} letter, as in \ar{سِينِمَا}, takes no sukoon. A \emph{glide} letter, the second element of a diphthong as in \ar{سَايْكُو}, takes sukoon. A \emph{rhotic} vowel is written as a mark followed by \ar{ر}, as in \ar{شَرْ}.

\section{Sukoon}
Sukoon marks a consonant followed by no vowel. In full density every such consonant carries it, including the last consonant of a word, as in \ar{لَاكْ} \emph{lack} and \ar{كَابِتَالِسْتْ} \emph{capitalist}. Sukoon records the absence of a vowel; it does not make a letter silent, and English double letters are not marked with shadda.

\section{Word-initial vowels}
A word that begins with a vowel takes a hamza seat: \ar{أ} for fatha and damma, \ar{إ} for kasra, as in \ar{أُفْ} \emph{of} and \ar{إِنْ} \emph{in}. An initial diphthong beginning with fatha and alef is written with \ar{آ}. A vowel that follows another vowel is hosted by a preceding \ar{ي} when there is one (\ar{ثِيُ}) and otherwise takes a seat.

\section{Frequent words}
Frequent function words have fixed spellings (Table~\ref{tab:common}). These come from the original fixed word table and from forms that the page and the typed passage share. They also resolve the contextual variation of \emph{the} and \emph{a}: the toolkit writes \ar{ذَا} before a vowel too, as the typed passage does, and offers the phonemic alternative \ar{ذِي} as an option.

\begin{table}[ht]\centering\small
\begin{tabular}{@{}ll@{\qquad}ll@{}}
\toprule
English & Spelling & English & Spelling\\
\midrule
the & \ar{ذَا} & is & \ar{إِزْ}\\
this & \ar{ذِسْ} & she & \ar{شِي}\\
that & \ar{ذَتْ} & it & \ar{إِتْ}\\
of & \ar{أُفْ} & its & \ar{إِتْسْ}\\
and & \ar{أَنْدْ} & as & \ar{أَزْ}\\
in & \ar{إِنْ} & into & \ar{إِنْتُو}\\
to & \ar{تُو} & with & \ar{وِذْ}\\
a & \ar{أَ} & are & \ar{أَرْ}\\
\bottomrule
\end{tabular}
\caption{Fixed spellings of frequent words.}\label{tab:common}
\end{table}

\section{Punctuation and numbers}
The comma, semicolon and question mark take their Arabic forms \ar{،}, \ar{؛}, \ar{؟}. A possessive apostrophe after $s$ is silent, as on the page. Digits default to Arabic-Indic numerals. A \code{spoken} option writes numbers as English number words in the notation (\emph{twenty five}), which avoids the earlier failure of rendering numbers as Arabic number words.

\section{Letter sound and assigned function}
Each Arabic letter has an ordinary sound in Arabic, and in this notation it also has an assigned function in writing English. The two often coincide, as with \ar{ب} for \emph{b} and \ar{ل} for \emph{l}, and sometimes they differ. The letter \ar{غ} is a voiced uvular fricative in Arabic, and it is used here for the hard English \emph{g} by convention. The same holds for \ar{ث} and \ar{ذ}, which are close to English \emph{th} and \emph{dh} sounds, and for \ar{ب}, which stands for both \emph{b} and \emph{p}. The letters \ar{و} and \ar{ي} serve as consonants (\emph{w} and \emph{y}) and also as vowel letters.

Knowing Arabic script is therefore a useful starting point, since the letter shapes, the direction of writing and the role of the marks are already familiar. It does not supply every reading rule. A reader who pronounces \ar{غ} as in Arabic will mispronounce \emph{good} and \emph{begin}, and a reader who expects \ar{ب} to be only \emph{b} will miss the \emph{p} in \emph{script}. The tables in this chapter give the assigned function, and the Arabic sound is only a mnemonic where it helps.

\section{Short vowels in practice}
Short vowels are written with a single mark on the preceding consonant and no vowel letter. The words below show the five English short vowels, each followed by one consonant or a cluster.

\begin{table}[ht]\centering\small
\begin{tabular}{@{}lll@{}}
\toprule
English & Arabic script & Respelling\\
\midrule
cat & \ar{كَاتْ} & \textit{kat}\\
bed & \ar{بِدْ} & \textit{bed}\\
sit & \ar{سِتْ} & \textit{sit}\\
hot & \ar{هُوتْ} & \textit{haat}\\
cut & \ar{كَتْ} & \textit{kut}\\
put & \ar{بُتْ} & \textit{put}\\
pen & \ar{بِنْ} & \textit{pen}\\
map & \ar{مَابْ} & \textit{map}\\
top & \ar{تُوبْ} & \textit{taap}\\
sun & \ar{سَنْ} & \textit{sun}\\
\bottomrule
\end{tabular}
\caption{Short vowels, written with a mark only.}\label{tab:wshort}
\end{table}

\section{Long vowels and the final e}
English marks long vowels in several ways: a doubled letter, a vowel digraph, or a final silent \emph{e}. The dictionary reading already gives the long vowel in each case, so the notation writes a vowel letter after the mark. The final \emph{e} itself is silent and is not written.

\begin{table}[ht]\centering\small
\begin{tabular}{@{}lll@{}}
\toprule
English & Arabic script & Respelling\\
\midrule
make & \ar{مَيْكْ} & \textit{meyk}\\
time & \ar{تَايْمْ} & \textit{taym}\\
home & \ar{هُومْ} & \textit{hohm}\\
use & \ar{يُوسْ} & \textit{yoos}\\
cake & \ar{كَيْكْ} & \textit{keyk}\\
like & \ar{لَايْكْ} & \textit{layk}\\
phone & \ar{فُونْ} & \textit{fohn}\\
tune & \ar{تُونْ} & \textit{toon}\\
\bottomrule
\end{tabular}
\caption{Long vowels and glides written for the silent-e pattern.}\label{tab:wlong}
\end{table}

\section{Doubled letters and vowel digraphs}
When the English spelling uses two vowel letters for one vowel, the notation keeps a waw or a yeh. This applies to the short vowel of \emph{book} as much as to the long vowel of \emph{food}. The spelling carries the weight here and the Arabic letter signals it. Chapter~\ref{ch:vowels} explains why \emph{book} is \ar{بُوكْ} and not \ar{بُكْ}.

\begin{table}[ht]\centering\small
\begin{tabular}{@{}lll@{}}
\toprule
English & Arabic script & Respelling\\
\midrule
book & \ar{بُوكْ} & \textit{buk}\\
look & \ar{لُوكْ} & \textit{luk}\\
good & \ar{غُودْ} & \textit{gud}\\
foot & \ar{فُوتْ} & \textit{fut}\\
wood & \ar{وُودْ} & \textit{wud}\\
food & \ar{فُودْ} & \textit{food}\\
moon & \ar{مُونْ} & \textit{moon}\\
school & \ar{سْكُولْ} & \textit{skool}\\
too & \ar{تُو} & \textit{too}\\
two & \ar{تُو} & \textit{too}\\
blue & \ar{بْلُو} & \textit{bloo}\\
true & \ar{تْرُو} & \textit{troo}\\
\bottomrule
\end{tabular}
\caption{Words with a written double o, and related spellings.}\label{tab:woo}
\end{table}

\begin{table}[ht]\centering\small
\begin{tabular}{@{}lll@{}}
\toprule
English & Arabic script & Respelling\\
\midrule
see & \ar{سِي} & \textit{see}\\
tree & \ar{تْرِي} & \textit{tree}\\
green & \ar{غْرِينْ} & \textit{green}\\
sleep & \ar{سْلِيبْ} & \textit{sleep}\\
head & \ar{هِدْ} & \textit{hed}\\
bread & \ar{بْرِدْ} & \textit{bred}\\
read & \ar{رِدْ} & \textit{red}\\
meet & \ar{مِيتْ} & \textit{meet}\\
meat & \ar{مِيتْ} & \textit{meet}\\
dream & \ar{دْرِيمْ} & \textit{dreem}\\
dead & \ar{دِدْ} & \textit{ded}\\
\bottomrule
\end{tabular}
\caption{Words with ee and ea.}\label{tab:wee}
\end{table}

\section{Glides and diphthongs}
The diphthongs of English are written as a vowel mark followed by a glide letter that takes sukoon, so \emph{day} is \ar{دَيْ}, \emph{light} is \ar{لَايْتْ} and \emph{voice} is \ar{فُويْسْ}. This is the placement seen on the handwritten page.

\begin{table}[ht]\centering\small
\begin{tabular}{@{}lll@{}}
\toprule
English & Arabic script & Respelling\\
\midrule
rain & \ar{رَيْنْ} & \textit{reyn}\\
day & \ar{دَيْ} & \textit{dey}\\
light & \ar{لَايْتْ} & \textit{layt}\\
night & \ar{نَايْتْ} & \textit{nayt}\\
voice & \ar{فُويْسْ} & \textit{voys}\\
flower & \ar{فْلَاوَرْ} & \textit{flower}\\
hour & \ar{آوَرْ} & \textit{ower}\\
boy & \ar{بُويْ} & \textit{boy}\\
toy & \ar{تُويْ} & \textit{toy}\\
noise & \ar{نُويْزْ} & \textit{noyz}\\
now & \ar{نَاوْ} & \textit{now}\\
how & \ar{هَاوْ} & \textit{how}\\
\bottomrule
\end{tabular}
\caption{Glides and diphthongs.}\label{tab:wai}
\end{table}

\section{R-coloured vowels}
A vowel followed by \emph{r} in the dictionary is written as a vowel mark and \ar{ر}. The mark of an unstressed final syllable is fatha in every attested form, as in \ar{نَيْتْشَر} and \ar{هَوِبَرْ}. Stressed syllables take their mark from the spelling.

\begin{table}[ht]\centering\small
\begin{tabular}{@{}lll@{}}
\toprule
English & Arabic script & Respelling\\
\midrule
car & \ar{كَارْ} & \textit{kaar}\\
far & \ar{فَارْ} & \textit{faar}\\
her & \ar{هِرْ} & \textit{her}\\
bird & \ar{بِرْدْ} & \textit{berd}\\
turn & \ar{تَرْنْ} & \textit{tern}\\
word & \ar{وُرْدْ} & \textit{werd}\\
door & \ar{دُورْ} & \textit{dawr}\\
more & \ar{مُورْ} & \textit{mawr}\\
fire & \ar{فَايَرْ} & \textit{fayer}\\
tire & \ar{تَايَرْ} & \textit{tayer}\\
\bottomrule
\end{tabular}
\caption{R-coloured vowels.}\label{tab:wr}
\end{table}

\section{Clusters and sukoon}
English has many consonant clusters that Arabic does not allow at the start of a word. The notation writes them with sukoon on all but the last consonant, so the cluster \emph{str} in \emph{street} appears as \ar{سْتْرِيتْ}. Word-final consonants also take sukoon in the full density. The \texttt{medial} density removes the final sukoon, which leaves clusters inside words intact.

\begin{table}[ht]\centering\small
\begin{tabular}{@{}lll@{}}
\toprule
English & Arabic script & Respelling\\
\midrule
string & \ar{سْتْرِنْغْ} & \textit{string}\\
spring & \ar{سْبْرِنْغْ} & \textit{spring}\\
splash & \ar{سْبْلَاشْ} & \textit{splash}\\
street & \ar{سْتْرِيتْ} & \textit{street}\\
strength & \ar{سْتْرِنْغْكْثْ} & \textit{strengkth}\\
twelfth & \ar{تْوِلْفْثْ} & \textit{twelfth}\\
asked & \ar{أَسْكْتْ} & \textit{askt}\\
texts & \ar{تِكْسْتْسْ} & \textit{teksts}\\
\bottomrule
\end{tabular}
\caption{Clusters, written with sukoon.}\label{tab:wclusters}
\end{table}
